%% kapa: exam=1 round=25 year=2014 period=1 form=A subject=economics subject_ko=경제원론 seq=36 number=76
%% source: Q-Net 감정평가사 기출문제 (cst00309) artlSeq=5101207 file=extracted/1008/5101207/2014년 제25회 감정평가사 1차시험 기출문제/1교시(A형).hwp
%% conversion: kordoc 4.24.51 -> tools/gen1.py (auto); review=proofread-sonnet; flags=none
\begin{question}{76}
신고전학파의 투자이론에 관한 설명으로 옳지 않은 것은? (단, 감가상각률과 자본재 가격의 변화율 및 조세의 영향은 고려하지 않는다.)
\begin{choices}
  \item 실질이자율이 상승하면 기업의 투자는 감소한다.
  \item 실질이자율이 하락하면 자본의 한계생산도 하락한다.
  \item 경제 전체의 기술진보로 인하여 자본의 한계생산이 높아지면 기업의 투자수요는 증가한다.
  \item 경제 전체의 기술진보로 인하여 자본의 한계생산이 높아지면 이자율은 상승한다.
  \item 감가상각률을 고려하지 않으므로 자본재 1단위에 대한 투자의 기회비용은 자본재 1단위의 매매가격과 같다.
\end{choices}
\end{question}
\answer{5}
